% Generated draft: requires mathematical and source review.
\documentclass[11pt]{article}
\usepackage[margin=25mm]{geometry}
\usepackage{fontspec}
\setmainfont{Latin Modern Roman}
\usepackage{amsmath,amssymb,mathtools,unicode-math}
\setmathfont{Latin Modern Math}
\usepackage{enumitem,needspace,xurl,hyperref}
\setlength{\parindent}{0pt}
\setlength{\parskip}{5pt}
\setlength{\emergencystretch}{3em}
\newcommand{\N}{\mathbb{N}}
\newcommand{\Z}{\mathbb{Z}}
\newcommand{\Q}{\mathbb{Q}}
\newcommand{\R}{\mathbb{R}}
\newcommand{\C}{\mathbb{C}}
\newcommand{\F}{\mathbb{F}}
\newcommand{\A}{\mathbb{A}}
\newcommand{\PP}{\mathbb{P}}
\newcommand{\E}{\mathbb{E}}
\newcommand{\eps}{\varepsilon}
\newcommand{\dd}{\,\mathrm{d}}
\begin{document}
\section*{UnsolvedMath ID: 11100044}
\textbf{Source problem code:} catalogue entry
\subsection{Definitions and mathematical statement}
A profinite group is a compact totally disconnected topological group, expressible as an inverse limit of finite groups. The source proposes \(\Z_p\), with its profinite topology, as a first example, and the Morava stabilizer group as the main example. The notation does not mean the additive field \(\Q_p\).

\emph{Equivariant-homotopy problem 2 (Mark Hovey).} Construct an equivariant stable homotopy theory for profinite structure groups, analogous in the intended sense to the compact-Lie-group theory. Explain how Morava \(E\)-theory becomes an object of the category, carrying the intended stabilizer action. The source does not specify the prime and height, ordinary versus extended stabilizer, continuity model, spectrum category, weak equivalences, or stabilization data.

A precise answer must distinguish spectra with homotopy-coherent or continuous group action from a genuine equivariant theory with the desired representation suspension, fixed-point, and transfer structures. Homotopy fixed-point constructions alone do not automatically identify all such theories. For Morava \(E_n\), an appropriate continuous/profinite action must be specified rather than treating the stabilizer as a discrete group or assuming a discrete-action model automatically captures its completed coefficients.

Hovey's opening assertion that only compact-Lie-group theory was then available is historical. The supplied registry marks the problem solved, and its summary points to profinite spectra, Galois descent and homotopy-fixed-point frameworks, including \(E_n^{hG}\) and higher-real-\(K\)-theory contexts. Its claim that the theory is ``fully developed'' does not specify whether it answers the genuine, Borel, continuous, or another intended version. Those supplied progress claims require matching precise definitions; no equivalence among all profinite equivariant frameworks is asserted.
\subsection{Short English statement}
Develop profinite equivariant stable homotopy theory, beginning with \(\Z_p\), and place Morava \(E\)-theory with its stabilizer action inside it. The registry reports later solutions via descent and homotopy fixed points, but the precise continuous and genuine-equivariant conventions remain unspecified.
\subsection{Sources}
\begin{itemize}
\item Local UnsolvedMath registry, permanent ID 11100044, source code catalogue entry; proposer Mark Hovey, ``Equivariant homotopy,'' item 2.
\item \emph{Hovey's Algebraic Topology Problem List}, surviving mirror: \url{https://www-users.cse.umn.edu/~tlawson/hovey/}. Recorded access 2026-07-29; extraction html-ordered-list. The mirror contains 65 top-level items across nine category pages.
\item The supplied summary names Behrens--Davis, ``The homotopy fixed point spectra of profinite Galois extensions,'' and unspecified Davis--Lawson and other profinite-spectra work, without full bibliography or exact theorem scopes. Its Morava \(E_n\), \(EO_n\), and open-subgroup discussion is truncated.
\item Supplied status solved and SOLVED-IN-LITERATURE. The summary describes profinite equivariant theory as fully developed through homotopy-fixed-point/descent work, a supplied assessment needing definition-specific review. Generic background retains NEEDS\_REVIEW for a modern primary-literature audit and redistribution terms. The registry assigns L4 whereas background records default L3.
\end{itemize}
\end{document}
